% Part: set-theory
% Chapter: ordinals
% Section: iso

\documentclass[../../../include/open-logic-section]{subfiles}

\begin{document}

\olfileid{sth}{ordinals}{iso}
\olsection{ক্রমসমরূপতা}

সুক্রমবিন্যাস কতটা \emph{শক্তিশালী} ধারণা, তা
বোঝাতে প্রথমে সুক্রমবিন্যাসগুলির তুলনার একটি
পদ্ধতি দিই।

\begin{defn}
একটি \emph{সুক্রমবিন্যাস} হলো এমন একটি যুগল
$\tuple{A, <}$, যাতে $<$ সম্পর্কটি $A$-কে সুক্রমে
বিন্যস্ত করে। দুটি সুক্রমবিন্যাস $\tuple{A, <}$ এবং
$\tuple{B, \lessdot}$ \emph{ক্রমসমরূপ} {যদি এবং কেবল যদি}
এমন একটি !!a{bijection} $f \colon A \to B$ থাকে, যাতে
$x < y$ যদি এবং কেবল যদি $f(x) \lessdot  f(y)$।
এ ক্ষেত্রে আমরা $\ordeq{\tuple{A, <}}{\tuple{B, \lessdot}}$
লিখি এবং $f$-কে একটি \emph{ক্রমসমরূপ চিত্রণ} বলি।
\end{defn}
\noindent
এর পর সংক্ষেপে আমরা ``ক্রমসমরূপ চিত্রণ''-এর বদলে
``সমরূপ চিত্রণ'' বলব। স্বজ্ঞাতভাবে, সমরূপ চিত্রণ হলো
গঠন অক্ষত রাখা !!{bijection}s। সমরূপ চিত্রণ সম্পর্কে
কয়েকটি সহজ তথ্য দেখি।

\begin{lem}\ollabel{isoscompose}
সমরূপ চিত্রণগুলির মিশ্রণও সমরূপ চিত্রণ। অর্থাৎ,
$f \colon A \to B$ এবং $g \colon B \to C$ সমরূপ
চিত্রণ হলে $(g \circ f) \colon A \to C$-ও সমরূপ
চিত্রণ।
\end{lem}
\begin{prob}
	\olref[sth][ordinals][iso]{isoscompose} প্রমাণ করুন।
\end{prob}
\begin{proof}
অনুশীলনের জন্য রাখা হলো।
\end{proof}
\begin{cor}\ollabel{ordisoisequiv}
	$\ordeq{X}{Y}$ একটি তুল্যতা সম্পর্ক।
\end{cor}
\begin{prop}\ollabel{ordisounique}
$\tuple{A, <}$ ও $\tuple{B, \lessdot}$ যদি
ক্রমসমরূপ সুক্রমবিন্যাস হয়, তবে তাদের মধ্যকার
সমরূপ চিত্রণটি একমাত্র।
\end{prop}

\begin{proof}
ধরি, $f$ ও $g$ দুটি সমরূপ চিত্রণ $A \to B$।
\olref[wo]{propwoinduction} অনুযায়ী আবেশ ব্যবহার
করে ফলটি প্রমাণ করব। একটি $a\in A$ নিই এবং
(আবেশের অনুমান হিসেবে) ধরি
$(\forall b < a)f(b) = g(b)$। এবার একটি $x \in B$ নিই।

যদি $x \lessdot f(a)$ হয়, তবে
$f^{-1}(x) < a$। তাই $f$ ও $g$ সমরূপ চিত্রণ
হওয়ায় $g(f^{-1}(x)) \lessdot g(a)$। কিন্তু
$f^{-1}(x) < a$ বলে আমাদের অনুমান থেকে
$x =f(f^{-1}(x)) = g(f^{-1}(x))$। সুতরাং
$x \lessdot g(a)$। একইভাবে, যদি $x \lessdot g(a)$,
তবে $x \lessdot f(a)$।

সার্বিকভাবে, $(\forall x \in B)(x \lessdot f(a) \liff x \lessdot
g(a))$। অতএব
\olref[sfr][rel][ord]{prop:extensionality-strictlinearorders}
অনুযায়ী $f(a) = g(a)$। তাই
\olref[wo]{propwoinduction} থেকে $(\forall
a \in A)f(a) = g(a)$।
\end{proof}\noindent
এ থেকে সুক্রমবিন্যাসের শক্তিমত্তা কিছুটা বোঝা যায়।
বিষয়টি আরও বুঝতে আরেকটু সংকেতলিপি দরকার।

\begin{defn}
$\tuple{A, <}$ একটি সুক্রমবিন্যাস এবং $a \in A$
হলে ধরি, $A_a = \Setabs{x \in A}{x < a}$।
$A_a$-কে $A$-এর যথার্থ \emph{প্রারম্ভিক খণ্ড}
বলি (এবং $A$-কেই $A$-এর অযথার্থ প্রারম্ভিক
খণ্ড হিসেবে মানি)। প্রারম্ভিক খণ্ডে $<$-এর সীমাবদ্ধন,
অর্থাৎ $\funrestrictionto{\mathord{<}}{A_a^2}$,
লিখি $<_a$।
\end{defn}
\noindent
এই সংকেতলিপিতে বলা ও প্রমাণ করা যায় যে,
কোনো সুক্রমবিন্যাস তার নিজের কোনো যথার্থ
প্রারম্ভিক খণ্ডের সঙ্গে সমরূপ নয়।

\begin{lem}\ollabel{wellordnotinitial}
$\tuple{A, <}$ সুক্রমবিন্যাস এবং $a \in A$ হলে
$\ordneq{\tuple{A, <}}{\tuple{A_a, <_a}}$।
\end{lem}

\begin{proof}
বিরোধের উদ্দেশ্যে ধরি, $f \colon A \to A_a$
একটি সমরূপ চিত্রণ। $f$ একটি একৈক সমাপতন এবং
$A_a \subsetneq A$ বলে,
\olref[wo]{wo:strictorder} ব্যবহার করে $A$-তে
$b \neq f(b)$ ধর্মবিশিষ্ট $<$-ক্ষুদ্রতম !!{element}
$b \in A$ নিই। আমরা দেখাব,
$(\forall x \in A)(x<b \liff x < f(b))$।
তাহলে
\olref[sfr][rel][ord]{prop:extensionality-strictlinearorders}
থেকে $b = f(b)$ পাওয়া যাবে; এটাই বিরোধ।

ধরি, $x < b$। $b$-এর নির্বাচনের জন্য $x = f(x)$।
আর $f$ সমরূপ চিত্রণ বলে $f(x) < f(b)$।
তাই $x < f(b)$।

ধরি, $x < f(b)$। $f$ সমরূপ চিত্রণ বলে
$f^{-1}(x) < b$; তাই $b$-এর নির্বাচনের জন্য
$f^{-1}(x) = x$। অতএব $x < b$।
\end{proof}

পরের ফলটি মোটামুটি বলে, সমরূপ চিত্রণের একটি
``প্রারম্ভিক খণ্ড''ও সমরূপ চিত্রণ:

\begin{lem}\ollabel{wellordinitialsegment}
ধরি, $\tuple{A, <}$ এবং $\tuple{B, \lessdot}$
সুক্রমবিন্যাস। যদি $f \colon A \to B$ সমরূপ
চিত্রণ এবং $a \in A$ হয়, তবে
$\funrestrictionto{f}{A_{a}} : A_a \to B_{f(a)}$
একটি সমরূপ চিত্রণ।
\end{lem}

\begin{proof}
$f$ সমরূপ চিত্রণ বলে:
	%$f$ সমরূপ চিত্রণ বলে $b < a$ যদি এবং কেবল যদি $f(b) \lessdot f(a)$; তাই
\begin{align*}
	\funimage{f}{A_a} &= \funimage{f}{\Setabs{x \in A}{x < a}}\\
	&= \funimage{f}{\Setabs{f^{-1}(y) \in A}{f^{-1}(y) < a}} \\
	&= \Setabs{y \in B}{y \lessdot f(a)} \\
	&=B_{f(a)}
\end{align*}
আর $\funrestrictionto{f}{A_a}$ ক্রম অক্ষত রাখে,
কারণ $f$-ও তা রাখে।
\end{proof}

পরের দুটি ফল থেকে জানা যাবে যে সুক্রমবিন্যাসগুলি
সব সময় তুলনীয়:

\begin{lem}\ollabel{lemordsegments}
ধরি, $\tuple{A, <}$ ও $\tuple{B, \lessdot}$
সুক্রমবিন্যাস। যদি
$\ordeq{\tuple{A_{a_1}, <_{a_1}}}{\tuple{B_{b_1}, \lessdot_{b_1}}}$
এবং
$\ordeq{\tuple{A_{{a_2}}, <_{a_2}}}{\tuple{B_{{b_2}},
\lessdot_{b_2}}}$ হয়, তবে ${a_1}  < {a_2} \text{ যদি এবং কেবল যদি }{b_1} \lessdot
{b_2}$।
\end{lem}

\begin{proof}
আমরা \emph{বাঁ থেকে ডানে} প্রমাণ করব; উল্টো দিকও
একই রকম। ধরি,
$\ordeq{\tuple{A_{a_1}, <_{a_1}}}{\tuple{B_{b_1},
\lessdot_{b_1}}}$ এবং
$\ordeq{\tuple{A_{{a_2}},
<_{a_2}}}{\tuple{B_{{b_2}}, \lessdot_{b_2}}}$
উভয়ই সত্য; এদের দ্বিতীয় সমরূপ চিত্রণটি
$f \colon A_{{a_2}} \to B_{{b_2}}$। ধরি,
${a_1} < {a_2}$। তাহলে
\olref{wellordinitialsegment} অনুযায়ী
$\ordeq{\tuple{A_{a_1}, <_{a_1}}}{\tuple{B_{f({a_1})},
\lessdot_{f({a_1})}}}$। সুতরাং
$\ordeq{\tuple{B_{b_1}, \lessdot_{b_1}}}{\tuple{B_{f({a_1})},
\lessdot_{f({a_1})}}}$। তাই
\olref{wellordnotinitial} থেকে
${b_1} = f({a_1})$। এখন
% BN-SRC-435: উৎসে চিত্রণের বিস্তৃতিকে সংজ্ঞাক্ষেত্র বলা; প্রকৃত বিস্তৃতি এখানে স্পষ্ট।
${b_1} \lessdot {b_2}$, কারণ f-এর বিস্তৃতি
$B_{{b_2}}$।
\end{proof}

\begin{thm}\ollabel{thm:woalwayscomparable}
যেকোনো দুটি সুক্রমবিন্যাসের একটি অন্যটির
কোনো একটি প্রারম্ভিক খণ্ডের সঙ্গে সমরূপ
(খণ্ডটি যথার্থ হওয়া জরুরি নয়)।
\end{thm}

\begin{proof}
ধরি, $\tuple{A, <}$ ও $\tuple{B, \lessdot}$
সুক্রমবিন্যাস। পৃথকীকরণ ব্যবহার করে ধরি,
\[
	f = \Setabs{\tuple{a, b} \in A \times B}{
		\ordeq{\tuple{A_a, <_a}}{\tuple{B_b, \lessdot_b}}}.
\]
\olref{lemordsegments} অনুযায়ী $f$-এর যেকোনো
$\tuple{a_1, b_1}, \tuple{a_2, b_2} \in f$-এর
জন্য $a_1 < a_2$ যদি এবং কেবল যদি
$b_1 \lessdot b_2$। তাই
$f \colon \dom{f} \to \ran{f}$ একটি সমরূপ
চিত্রণ।

যদি $a_2 \in \dom{f}$ এবং $a_1 < a_2$, তবে
\olref{wellordinitialsegment} অনুযায়ী
$a_1 \in \dom{f}$। অতএব $\dom{f}$ হলো
$A$-এর একটি প্রারম্ভিক খণ্ড। একইভাবে,
$\ran{f}$ হলো $B$-এর একটি প্রারম্ভিক খণ্ড।
বিরোধের উদ্দেশ্যে ধরি, উভয় খণ্ডই
\emph{যথার্থ}। তাহলে $A \setminus \dom{f}$-এর
$<$-ক্ষুদ্রতম !!{element} $a$ নিই; ফলে
$\dom{f} = A_a$। আর $B \setminus \ran{f}$-এর
$\lessdot$-ক্ষুদ্রতম !!{element} $b$ নিই; ফলে
$\ran{f} = B_b$। তাই $f \colon A_a \to B_b$
একটি সমরূপ চিত্রণ; সুতরাং
$\tuple{a, b} \in f$—বিরোধ।
\end{proof}

\end{document}
